\textbf{COLUNI 2023 -- Questão 03}\par\medskip Considere o trecho a seguir:

\par

“Quando você estava com alguém de quem eu não gostava, eu aumentava meu volume — e volume, você sabe, é coisa que não me falta — até chegar a níveis \underline{insuportáveis}, uma avalanche de decibéis.”

\par

SCLIAR, Moacyr. O rádio apaixonado. In: SCLIAR, Moacyr. \textbf{Histórias que os jornais não contam}. 3. ed. Porto Alegre: L\&M Editores, 2018. p.14.

\par

A alternativa em que a palavra sublinhada NÃO tem o mesmo processo de formação que o termo insuportáveis, em destaque no trecho acima é:\begin{enumerate}
\item “[...] Você não pode imaginar o sofrimento que isto me causa, mesmo porque você provavelmente acha
       que rádios são objetos \underline{inanimados}, sem vida própria [...]”. (SCLIAR, Moacyr. O rádio apaixonado. In: SCLIAR,
      Moacyr. \textbf{Histórias que os jornais não contam}. 3. ed. Porto Alegre: L\&M Editores, 2018. p.14)
\item “[...] Senti um verdadeiro choque, eu que já deveria estar \underline{acostumado} com eletricidade. Você fez de mim
       um ser vivo [...]”. (SCLIAR, Moacyr. O rádio apaixonado. In: SCLIAR, Moacyr. \textbf{Histórias que os jornais não contam}. 3. ed. Porto Alegre: L\&M Editores, 2018. p.14-15).
\item “[...] Num certo sentido você está certa: estou desconfigurado, estou \underline{desfigurado}, estou perturbado —
       mas tudo isso por causa do sofrimento que você me causou [...]”. (SCLIAR, Moacyr. O rádio apaixonado. In:
      SCLIAR, Moacyr. \textbf{Histórias que os jornais não contam}. 3. ed. Porto Alegre: L\&M Editores, 2018. p.15-16.)
\item “[...] Minha última mensagem é esta: nunca brinque com os sentimentos de um rádio apaixonado. Você
       vai ter, no mínimo, surpresas \underline{desagradáveis} [...].” (SCLIAR, Moacyr. O rádio apaixonado. In: SCLIAR, Moacyr.
      \textbf{Histórias que os jornais não contam}. 3. ed. Porto Alegre: L\&M Editores, 2018. p.15-16).\end{enumerate}
